%-----------------------------------------------------------------------------%
\chapter{PENDAHULUAN}\label{bab:1}
%-----------------------------------------------------------------------------%

%-----------------------------------------------------------------------------%
\section{Latar Belakang}\label{sec:latar-belakang}
%-----------------------------------------------------------------------------%

Peraturan perundang-undangan adalah peraturan tertulis yang memuat norma hukum yang mengikat secara
umum dan dibentuk atau ditetapkan oleh lembaga negara atau pejabat yang berwenang melalui prosedur
yang ditetapkan dalam peraturan perundang-undangan, sesuai yang dijelaskan dalam Undang-Undang
Republik Indonesia Nomor 12 Tahun 2011 (selanjutnya disebut dengan UU 12/2011). Peraturan
perundang-undangan dapat digunakan untuk menjawab pertanyaan-pertanyaan yang berkaitan dengan hukum,
seperti:

\begin{itemize}
  \item Peraturan apa saja yang berlaku pada suatu daerah?
  \item Apa hubungan suatu peraturan dengan peraturan lain?
  \item Apa saja peraturan yang mengatur suatu topik?
  \item Apa isi dari Pasal X ayat (Y) dari suatu peraturan perundang-undangan?
\end{itemize}

Pertanyaan-pertanyaan tersebut umumnya biasa dijawab oleh seorang ahli hukum, yang artinya,
penerapan peraturan perundang-undangan ini hanya dapat dilakukan dalam skala kecil dan biaya relatif
mahal. Komputer dapat menjadi alternatif untuk aplikasi tersebut dalam skala lebih besar dan biaya
lebih murah, jika peraturan perundang-undangan berupa data terstruktur yang dapat diolah oleh
komputer. Akan tetapi, data peraturan perundang-undangan umumnya dibuat dalam bentuk data
semi-terstruktur, yaitu berupa dokumen PDF dengan struktur konten seperti aturan penomoran dan
aturan pembagian (\textit{sectioning}), sementara penulisan isi peraturannya sendiri dalam bentuk
data tidak terstruktur yaitu teks bahasa manusia.

\begin{figure}[H]
  \centering
  \includegraphics[width=\textwidth]{pictures/terstruktur.png}
  \caption{Cuplikan ekstraksi data terstruktur dan data tidak terstruktur dari dokumen PDF UU 11/2020}
  \label{fig:ekstraksi-dokumen}
\end{figure}

\pic~\ref{fig:ekstraksi-dokumen} merupakan salah satu contoh ekstraksi informasi dari UU 11/2020
tentang Cipta Kerja. Pada gambar tersebut dapat dilihat bahwa sebagian informasi dapat diekstraksi
dari struktur dokumen peraturan perundang-undangan tak terstruktur yang berupa teks mentah (ditandai
dengan warna merah), dan sebagian informasi lainnya dapat diekstraksi dari data  (ditandai dengan
warna biru). Pada gambar tersebut, dalam mengekstraksi informasi ``UU 11/2020 Pasal 6 merujuk kepada
UU 11/2020 Pasal 5 ayat (1) huruf a'', manusia pertama-tama perlu membaca dari awal teks UU serta
pasal tersebut sampai menemukan pola \textit{string} terkait yang merujuk kepada suatu komponen
dokumen peraturan perundang-undangan, dan kemudian meng-\textit{infer} informasi lengkap berdasarkan
konteks komponen tersebut. Dengan kata lain, proses pencarian dan konversi \textit{string} ``Pasal 5
ayat (1) huruf a'' menjadi informasi lengkap ``UU No.11 Tahun 2020 Pasal 5 ayat (1) huruf a'' perlu
dilakukan secara manual oleh manusia.

\begin{figure}[H]
  \centering
  \includegraphics{pictures/outline.png}
  \caption{Bagan Alir Pembuatan dan Query KG Peraturan Perundang-undangan}
  \label{fig:konstruksi-dan-query}
\end{figure}

\textit{Knowledge graph} adalah struktur data yang (i) menggambarkan entitas dunia nyata dan
keterkaitannya, terorganisir dalam grafik, (ii) mendefinisikan kemungkinan kelas dan hubungan
entitas dalam skema, (iii) memungkinkan entitas memiliki keterkaitan satu sama lain dan (iv)
mencakup berbagai topik domain \citep{paulheim_2016}. Dengan melakukan ekstraksi data dari data
terstruktur dan tidak terstruktur dari dokumen, peraturan perundang-undangan dapat diubah menjadi
\textit{knowledge graph}.  Sebagai manfaatnya, data peraturan perundang-undangan yang sudah
dikonvesi menjadi \textit{knowledge graph} kemudian dapat dilakukan fitur-fitur berbasis
\textit{query} seperti \textit{question answering}, \textit{reasoning} dan visualisasi data.
\pic~\ref{fig:konstruksi-dan-query} menampilkan proses konstruksi \textit{knowledge graph} dari
beberapa berkas PDF peraturan perundang-undangan dan \textit{query} pada KG tersebut. Pada skripsi
ini, akan dilakukan konversi dokumen peraturan perundang-undangan menjadi data terstruktur dalam
bentuk \textit{knowledge graph}, dan memberikan contoh pengaplikasian dari \textit{knowledge graph}
peraturan perundang-undangan tersebut.

Beberapa usaha telah dilakukan untuk membuat \textit{vocabulary} untuk \textit{knowledge graph}
peraturan perundang-undangan seperti ELI (\url{https://eur-lex.europa.eu/eli-register/about.html}),
Schema.org (\url{https://schema.org/Legislation}), dan Wikidata
(\url{https://www.wikidata.org/wiki/Wikidata:WikiProject_Law}).


\begin{listing}[H]
  \begin{minted}[frame=single, breaklines, fontsize=\footnotesize]{turtle}
@prefix schema: <http://schema.org/> .
@prefix xsd: <http://www.w3.org/2001/XMLSchema#> .

<http://data.europa.eu/eli/dir/1980/181/oj>
  a schema:Legislation ;
  schema:about <http://eurovoc.europa.eu/1896> ;
  schema:isBasedOn <http://njh.me/11957E100.html> ;
  schema:legislationDate "1979-12-20"^^schema:Date ;
  schema:legislationPassedBy "Council of the European Union"^^xsd:string ;
  schema:name "Council Directive 80/181/EEC of 20 December 1979 on the approximation of the laws of the Member States relating to units of measurement and on the repeal of Directive 71/354/EEC"^^xsd:string ;

<http://eurovoc.europa.eu/1896>
  a schema:Thing ;
  schema:name "Metrology"^^xsd:string .

<http://njh.me/11957E100.html> a schema:Legislation .
  \end{minted}
  \caption{Contoh \textit{legal knowledge graph} dengan \textit{vocabulary} Schema.org}
  \label{listing:schema-org-example}
\end{listing}

Pada \lst~\ref{listing:schema-org-example} diberikan contoh \textit{knowledge graph} yang dibuat
dengan \textit{vocabulary} Schema.org. \textit{Knowledge graph} memuat metadata peraturan seperti
tentang apa peraturan tersebut (\code{schema:about}), apa dasar dari peraturan tersebut
(\code{schema:isBasedOn}), dan nama dari peraturan tersebut (\code{schema:name}).  Akan tetapi
\textit{vocabulary-vocabulary} yang disebutkan diatas tidak dapat merepresentasikan struktur
peraturan. Selain itu, \textit{Vocabulary} yang disebutkan diatas dirancang untuk digunakan secara
internasional atau khusus \textit{European Union} seperti ELI, sehingga tidak dapat
merepresentasikan data yang hanya terdapat pada peraturan perundang-undangan di Indonesia. Sehingga
dengan membuat \textit{vocabulary} untuk struktur data dan peraturan perundang-undangan di
Indonesia, diharapkan dapat menjadi kontribusi dalam mendukung pembuatan \textit{knowledge graph}
untuk dokumen peraturan perundang-undangan.

%-----------------------------------------------------------------------------%
\section{Permasalahan}\label{sec:permasalahan}
%-----------------------------------------------------------------------------%

Berikut ini adalah rumusan permasalahan dari penelitian yang dilakukan:
\begin{itemize}
  \itemsep1pt\parskip0pt\parsep0pt
  \item Bagaimana cara membuat \textit{legal ontology} untuk memperjelas \textit{requirement} dan
        mempermudah konversi dokumen PDF menjadi \textit{knowledge graph}?
  \item Bagaimana konversi dokumen peraturan perundang-undangan menjadi \textit{knowledge graph} dapat
        dilakukan secara otomatis?
  \item Apa saja contoh aplikasi dari knowledge graph peraturan perundang-undangan tersebut?
\end{itemize}

% // TODO: bagaimana melakukan proof of concept untuk yang udah dibuat, detail untuk maintain, general
% large scale, evaluasi konversi

% // TODO: bolongnya State of the art adalah ...

% // TODO: 1. Bagaimana merancang suatu sistem untuk melakukan automated population? 
% 2. memperjelas requirement dan menangkap informasi yang terkandung di dalam dokumen legal yang ada, sehingga mendukung konversi pdf menjadi knowledge graph.
% 3. bagaimana evaluasi use case terhadap kg peraturan perundang-undangan tersebut serta evaluasi konversi / populasi dalam skala besar

%-----------------------------------------------------------------------------%
\section{Batasan Permasalahan}\label{sec:batasan-permasalahan}
%-----------------------------------------------------------------------------%

% // TODO: fokus yang mainteined (UU + non UU), dan coba parsing UU hasil scraping

% // TODO: batasan: tipe query/question nya belum yang super complex (contoh: jika saya usia 12 tahun, apakah boleh kerja menurut UU ketenaga kerjaan?)

%-----------------------------------------------------------------------------%
\section{Sistematika Penulisan}\label{sec:sistematika-penulisan}
%-----------------------------------------------------------------------------%